% !TEX TS-program = pdflatex
% !TeX program = pdflatex
% !TEX encoding = UTF-8
% !TEX spellcheck = fr

\documentclass[xcolor=table]{beamer}


%\usepackage{fullpage}
%\usepackage[left=2.8cm,right=2.2cm,top=2 cm,bottom=2 cm]{geometry}
\setbeamersize{text margin left=10pt,text margin right=10pt}


\usepackage{amsmath,amssymb} 
\usepackage{fontenc}
\usepackage{textcomp}
\usepackage[utf8]{inputenc}
\usepackage{CJKutf8}
\usepackage[french]{babel}
\usepackage{arabtex,acolor}
\usepackage{txfonts}
\usepackage{tipa}
\usepackage[]{graphicx}
\usepackage{multirow}
\usepackage{hyperref}
\usepackage{colortbl}
\usepackage{tabularray}
\usepackage{listingsutf8}
\usepackage{wrapfig}
\usepackage{multicol}
\usepackage[export]{adjustbox} %for images in table, also for frame
\usepackage[many]{tcolorbox}
\usepackage{wasysym}
\usepackage[french,lined]{algorithm2e}
\usepackage{alltt} %verbatim with commands
\usepackage{longtable}
\usepackage{tabu}


\definecolor{lightblue}{HTML}{D0D2FF}
\definecolor{lightyellow}{HTML}{FFFFAA}
\definecolor{darkblue}{HTML}{0000BB}
\definecolor{olivegreen}{HTML}{006600}
\definecolor{darkgreen}{HTML}{008B45} %009B55
\definecolor{violet}{HTML}{6600CC}
%\definecolor{deeppink}{HTML}{FF1493}
\definecolor{orangey}{HTML}{FFBB00}

\usetheme{Warsaw} % Antibes Boadilla Warsaw

\beamertemplatenavigationsymbolsempty


\let\oldcite\cite
\renewcommand{\cite}[1]{{\bfseries\color{orangey}\oldcite{#1}}}

\tcbuselibrary{listings}

\newcommand{\kurl}[1]{{\scriptsize\bfseries\color{orangey}\url{#1}}}


%\renewcommand{\baselinestretch}{1.5}

\def\supit#1{\raisebox{0.8ex}{\small\it #1}\hspace{0.05em}}

\AtBeginSection{%
	\begin{frame}
		\sectionpage
	\end{frame}
}

\newcommand{\rottext}[2]{%
	\rotatebox{90}{%
	\begin{minipage}{#1}%
		\raggedleft#2%
	\end{minipage}%
	}%
}


\institute{ %
Laboratoire de la Communication dans les Systèmes Informatiques (LCSI)

\vspace{6pt}
École  nationale Supérieure d'Informatique (ESI, ex. INI), Alger, Algérie
}
\author[ \textbf{\footnotesize\insertframenumber/\inserttotalframenumber} \hspace*{1cm} ESI - ARIES Abdelkrime (Master 2022/2023)] %
{ARIES Abdelkrime}
%\titlegraphic{\includegraphics[height=1cm]{../img/esi-logo.png}%\hspace*{4.75cm}~
	
	
	\date{Master 2022/2023} %\today

\titlegraphic{%
	\includegraphics[height=1cm]{../img/esi-logo.png}%
	\hspace{2cm}
	\includegraphics[height=1cm]{../img/lcsi-logo.png}%
}


%\setbeamertemplate{headline}{}



\newcommand{\keyword}[1]{\textcolor{red}{\bfseries\itshape #1}}
\newcommand{\expword}[1]{\textcolor{olivegreen}{#1}}
\newcommand{\optword}[1]{\textcolor{violet}{\bfseries #1}}

\makeatletter
\newcommand\mysphere{%
	\parbox[t]{10pt}{\raisebox{0.2pt}{\beamer@usesphere{item projected}{bigsphere}}}}
\makeatother

%\let\oldtabular\tabular
%\let\endoldtabular\endtabular
%\renewenvironment{tabular}{\rowcolors{2}{white}{lightblue}\oldtabular\rowcolor{blue}}{\endoldtabular}


\NoAutoSpacing %french autospacing after ":"

\def\graphpath{}

\newcommand{\changegraphpath}[1]{\def\graphpath{#1}}


\newcommand{\vgraphpage}[2][.84\textheight]{%
%	\begin{center}%
		\includegraphics[height=#1]{\graphpath #2}%
%	\end{center}%
}

\newcommand{\hgraphpage}[2][\textwidth]{%
%	\begin{center}%
		\includegraphics[width=#1]{\graphpath #2}%
%	\end{center}%
}

\newcommand{\graphpage}[2][]{%
	\includegraphics[#1]{\graphpath #2}%
}

\bibliographystyle{apalike}

\newcommand{\insertbibliography}[2]{
	\appendix
	\section*{Bibliographie}
	\nocite{#2}
%	\makeatletter % to change template
%	\setbeamertemplate{headline}[default] % not mandatory, but I though it was better to set it blank
%	\def\beamer@entrycode{\vspace*{-\headheight}} % here is the part we are interested in :)
%	\makeatother
	\begin{multicols*}{2}[\frametitle{\insertsection} \usebeamertemplate*{frametitle}]%\usebeamertemplate*{frametitle}\frametitle{Références}
		\tiny
		\bibliography{#1}
	\end{multicols*}
}

\definecolor{my-grey}{RGB}{233, 233, 233}

\newcommand{\insertlicence}{
	\begin{frame}[plain]
	\frametitle{Licence : CC-BY 4.0}
%	\framesubtitle{Licence: CC-BY-NC 4.0}

	\begin{tcolorbox}[colback=cyan,
		colframe=cyan,  
		arc=0pt,outer arc=0pt,
		valign=top, 
		halign=center,
		width=\textwidth]
		
		\includegraphics[width=.5cm]{../img/licence/cc_icon_white_x2.png}
		\includegraphics[width=.5cm]{../img/licence/attribution_icon_white_x2.png}
		
		\color{white}
		\bfseries Attribution 4.0 International (CC BY 4.0) \\
		\tiny \url{https://creativecommons.org/licenses/by/4.0/deed.fr}
		
	\end{tcolorbox}\vspace{-.5cm}
	\begin{tcolorbox}[colback=my-grey,
		colframe=my-grey,  
		center, arc=0pt,outer arc=0pt,
		valign=top, 
		halign=left,
		width=\textwidth]
		
		\tiny
		
		\begin{center}
			\bfseries\Large
			Vous êtes autorisé à :
		\end{center}
		
		\begin{minipage}{0.83\textwidth}
			\begin{itemize}
				\item[] \textbf{Partager} — copier, distribuer et communiquer le matériel par tous moyens et sous tous formats
				\item[] \textbf{Adapter} — remixer, transformer et créer à partir du matériel
				pour toute utilisation, y compris commerciale.
			\end{itemize}
		\end{minipage}
		\begin{minipage}{0.15\textwidth}
			\includegraphics[width=\textwidth]{../img/licence/FreeCulturalWorks_seal_x2.jpg}
		\end{minipage}
	
		
		\begin{center}
			\bfseries\Large
			Selon les conditions suivantes :
		\end{center}
		
		\begin{itemize}
			\item[] \textbf{Attribution} — Vous devez créditer l'Œuvre, intégrer un lien vers la licence et indiquer si des modifications ont été effectuées à l'Oeuvre. Vous devez indiquer ces informations par tous les moyens raisonnables, sans toutefois suggérer que l'Offrant vous soutient ou soutient la façon dont vous avez utilisé son Oeuvre. 
			\item[] \textbf{Pas de restrictions complémentaires} — Vous n'êtes pas autorisé à appliquer des conditions légales ou des mesures techniques qui restreindraient légalement autrui à utiliser l'Oeuvre dans les conditions décrites par la licence.
		\end{itemize}
		
	\end{tcolorbox}
	
%	\begin{center}
%		\bfseries Attribution 4.0 International (CC BY 4.0)
%		\url{https://creativecommons.org/licenses/by/4.0/deed.fr}
%	\end{center}

%	\tiny
%
%	Vous êtes autorisé à : 
%	\begin{itemize}
%		\item \textbf{Partager} — copier, distribuer et communiquer le matériel par tous moyens et sous tous formats
%		\item \textbf{Adapter} — remixer, transformer et créer à partir du matériel
%	\end{itemize}
%	
%	Selon les conditions suivantes : 
%	\begin{itemize}
%		\item \textbf{Attribution} — Vous devez créditer l'Œuvre, intégrer un lien vers la licence et indiquer si des modifications ont été effectuées à l'Oeuvre. Vous devez indiquer ces informations par tous les moyens raisonnables, sans toutefois suggérer que l'Offrant vous soutient ou soutient la façon dont vous avez utilisé son Oeuvre.
%		\item \textbf{Pas d'Utilisation Commerciale} — Vous n'êtes pas autorisé à faire un usage commercial de cette Oeuvre, tout ou partie du matériel la composant. 
%		\item \textbf{Pas de restrictions complémentaires} — Vous n'êtes pas autorisé à appliquer des conditions légales ou des mesures techniques qui restreindraient légalement autrui à utiliser l'Oeuvre dans les conditions décrites par la licence.
%	\end{itemize}

	\end{frame}
}

\settowidth{\leftmargini}{\usebeamertemplate{itemize item}}
\addtolength{\leftmargini}{\labelsep}

\AtBeginDocument{
	\newcolumntype{L}[2]{>{\vbox to #2\bgroup\vfill\flushleft}p{#1}<{\egroup}} 
	
	\begin{frame}[plain]
		\maketitle
	\end{frame}

	\insertlicence
}


% needs etoolbox; to break links after -
\appto\UrlBreaks{\do\-}


\makeatletter
\newcommand{\xRightarrow}[2][]{\ext@arrow 0359\Rightarrowfill@{#1}{#2}}
\makeatother


\usefonttheme{structurebold}
%\usefonttheme{professionalfonts}

\title[TALN : 10- Quelques applications]%
{Traitement automatique du langage naturel\\Chapitre 10 : Quelques applications} 

\changegraphpath{../img/app/}

\begin{document}
	
\begin{frame}
\frametitle{Traitement automatique du langage naturel}
\framesubtitle{Quelques applications : Introduction}

%\begin{exampleblock}{Exemple d'un texte en français}
%	\begin{center}
%		\Large\bfseries
%		L'hiver est un des quatre saisons de l'année. 
%		Le rapport était bien exposé. 
%		La pluie est une des caractéristiques de cette saison.
%		Le sujet était la saison de l'hiver.
%	\end{center}
%\end{exampleblock}
%
%\begin{itemize}
%	\item Est-ce que vous pouvez comprendre de quoi s'agit-il ?
%	\item S'il y a un problème, lequel ?
%	\item Est-ce qu'on peut améliorer ce texte ?
%\end{itemize}

\end{frame}

%\begin{frame}
%\frametitle{Traitement automatique du langage naturel}
%\framesubtitle{Cohérence du discours :  Un peu d'humour}
%
%\begin{center}
%	\vgraphpage{humour1.jpg}
%\end{center}
%
%\end{frame}

\begin{frame}
\frametitle{Traitement automatique du langage naturel}
\framesubtitle{Quelques applications : Plan}

\begin{multicols}{2}
%	\small
\tableofcontents
\end{multicols}
\end{frame}

%===================================================================================
\section{Transformation}
%===================================================================================

\begin{frame}
	\frametitle{Quelques applications}
	\framesubtitle{Transformation}
	\begin{itemize}
		\item 
	\end{itemize}
\end{frame}

\subsection{Traduction automatique}

\begin{frame}
	\frametitle{Quelques applications : Transformation}
	\framesubtitle{Traduction automatique}
	\begin{itemize}
		\item 
	\end{itemize}
\end{frame}

\subsection{Résumé automatique}

\begin{frame}
	\frametitle{Quelques applications : Transformation}
	\framesubtitle{Résumé automatique : Classification d'un résumé automatique}
	\hgraphpage{sum-classif.pdf}
\end{frame}

\begin{frame}
	\frametitle{Quelques applications : Transformation}
	\framesubtitle{Résumé automatique : Approches du résumé automatique}
	\begin{columns}
		\begin{column}{0.32\textwidth}
			\begin{block}{\scriptsize\bfseries\cite{12-nenkova-mckeown}}
				\begin{itemize}
					\item représentation du sujet 
					:
					mots du sujet,
					fréquences, 
					analyse sémantique latente, 
					modèles de sujets bayésiens,
					clustering
					%			\begin{itemize}
					%				\item mots du sujet 
					%				\item fréquences
					%				\item analyse sémantique latente
					%				\item modèles de sujets bayésiens
					%			\end{itemize}
					\item représentation des indicateurs : 
					par graphes, 
					apprentissage automatique
				\end{itemize}
			\end{block}
		\end{column}
		\begin{column}{0.3\textwidth}
			\begin{block}{\scriptsize\bfseries\cite{12-lloret-palomar}}
				\begin{itemize}
					\item statistique 
					\item par graphes
					\item basée discours
					\item par apprentissage automatique
				\end{itemize}
			\end{block}
		\end{column}
		\begin{column}{0.28\textwidth}
			\begin{block}{\scriptsize\bfseries\cite{19-aries-al}}
				\begin{itemize}
					\item statistique 
					\item par graphes
					\item linguistique 
					\item par apprentissage automatique
				\end{itemize}
			\end{block}
		\end{column}
	\end{columns}
\end{frame}

\begin{frame}
	\frametitle{Quelques applications : Transformation}
	\framesubtitle{Résumé automatique : Approche statistique}
	
	\begin{itemize}
		\item \optword{Fréquence des mots} 
		
		\hspace{.5cm}Ex. $Score_\text{TF-IDF}(s_i) = \sqrt{\sum\limits_{w_{ik} \in s_i} (\text{TF-IDF}(w_{ik}))^2}$
		
		\item \optword{Position des phrases (ou des mots)}
		
		\hspace{.5cm}Ex. $ Score_\text{pos}(s_i) = \max (\frac{1}{i}, \frac{1}{|D| - i + 1}) $
		
		\item \optword{Taille (longueur) des phrases}
		
		\hspace{.5cm}Ex. $ Score_\text{taille}(s_i) = \left\lbrace 
		\begin{array}{lll}
		0 & si & (L_i \geq L_{min}) \\
		\frac{L_i - L_{min}}{L_{min}} & sinon & \\
		\end{array}
		\right. $
		
		\item \optword{Mots du titre et des sous-titres}
		
		\hspace{.5cm}Ex. $ Score_{titre}(s_i) = \frac{\sum_{e \in T \bigcap s_i}{\frac{tf(e)}{tf(e)+1}}}
		{\sum_{e \in T}{\frac{tf(e)}{tf(e)+1}}} $
		
		\item \optword{Centroid}, \optword{Frequent itemsets}, \optword{Analyse sémantique latente}, ...
	\end{itemize}
\end{frame}

\begin{frame}
	\frametitle{Quelques applications : Transformation}
	\framesubtitle{Résumé automatique : Approche statistique (Exemple : \cite{13-aries-al})}
	
	\hgraphpage{tcc-arch.pdf}
	
\end{frame}

\begin{frame}
	\frametitle{Quelques applications : Transformation}
	\framesubtitle{Résumé automatique : Approche statistique (Exemple: \cite{13-aries-al} - suite)}
	
	\begin{itemize}
		\item Un texte peut contenir plusieurs sujets, et une phrase peut discuter plusieurs sujets
		\begin{itemize}
			\item Regroupement de phrases avec similarité et seuil de regroupement ($Th$)
		\end{itemize}
		\item Une phrase est importante si elle peut représenter le maximum des sujets(cluters)
		\[ Score(s_i , c_j , f_k ) = 1 + \sum_{\phi \in s_i} {P(f_k=\phi | s_i \in c_j)} \]
		\[ Score(s_i , \bigcap_{j} c_j , F) =  %\propto 
		\prod_{j} \prod_{k} Score(s_i , c_j , f_k ) \]
		$ s $ : phrase, $ c $ : cluster, $ f $ : caractéristique, $ F $ : ensemble de caractéristiques, $ \phi $: observation de $ f $.
		\item $f$ : TF (Uni, Bi); Pos (intervalle de 10); Len (réel, pré-traité)
	\end{itemize}
	
\end{frame}

\begin{frame}
	\frametitle{Quelques applications : Transformation}
	\framesubtitle{Résumé automatique : Approche statistique (Exemple: \cite{15-oufaida-al})}
	\vspace{-6pt}
	\begin{itemize}
		\item Représentation des mots : embeddings pré-entrainés (Polyglot)
		\item Clustering pour extraire les sous-sujets du texte
		\begin{itemize}
			\item On peut trouver le mot $w_j \in S_2$ le plus similaire à un mot $w_i$
			
			$Match(w_i | S_2) = \arg\max_{w_j \in S_2} sim(Rep(w_i), Rep(w_j))$
			
			\item Lien entre deux phrases $S_2$ et $S_2$
			
			$Sim(S_1, S_2) = \frac{\sum_{w_i \in S_1} Match(w_i | S_2) + \sum_{w_j \in S_2} Match(w_j | S_1)}{|S_1| + |S_2|}$
		\end{itemize}
	
		\item Score des termes en utilisant mRMR
		\begin{itemize}
			\item Information mutuelle  
			$I(X, Y) = \sum\limits_{x \in X} \sum\limits_{y \in Y} p(x, y) \log \frac{p(x, y)}{p(x) p(y)}$
			\item Redondance des termes 
			$Redondance(T_i) = \frac{1}{|S|} \sum\limits_{j \in S} I(T_i, T_j)$
			
			\item Score final d'un terme : Mutual Information (Difference/Quotient)
			
			$MID \equiv \max_{t \in T} Pertinence(t) - Redondance(t)$

			$MIQ \equiv \max_{t \in T} Pertinence(t) / Redondance(t)$
		\end{itemize}
	
		\item Score d'un phrase à base des scores des termes mRMR (score pondéré par des poids)
	
		\item Lors de la sélection d'une phrase, on décrémente les poids des termes sélectionnés
	\end{itemize}
	
\end{frame}

\begin{frame}
	\frametitle{Quelques applications : Transformation}
	\framesubtitle{Résumé automatique : Approche par graphes}
	
	\begin{itemize}
		\item \optword{Propriétés du graphe}
		\begin{itemize}
			\item Bushy paths 
			
			\hspace{.5cm}$Score_{\#arcs}(s_i) = |\{ s_j : a(s_i, s_j) \in A / s_j \in S, s_i \neq s_j \}|$
			
			\item Aggregate Similarity
			
			\hspace{.5cm}$Score_{aggregate}(s_i) = \sum\limits_{(s_i, s_j) \in E} sim(s_i, s_j)$
		\end{itemize}
		\item \optword{Méthodes itératives}
		\begin{itemize}
			\item Mettre à jour les scores des nœuds à base des voisins 
			\item L'arrêt : un état d'équilibre (on peut plus mettre à jour)
			\item Ex. TextRank \cite{04-mihalcea-tarau}
			
			$WS(V_i) = ( 1 - d) + d * \sum\limits_{V_j \in In(V_i)} \frac{w_{ji}}{\sum\limits_{V_k \in Out(V_j)} w_{jk}} WS(V_j)$
			
			$w_{ij} = \frac{|\{w_k \text{ / } w_k \in S_i \text{ and } w_k \in S_j\}|}{\log(|S_i|) + \log(|S_j|)}$
			
			$ d $ : damping factor (en général, environ $ 0.85 $)
		\end{itemize}
	\end{itemize}

\end{frame}

\begin{frame}
	\frametitle{Quelques applications : Transformation}
	\framesubtitle{Résumé automatique : Approche par graphes (Exemple: \cite{21-aries-al})}
	
	\begin{center}
		\vgraphpage{gc-archi.pdf}
	\end{center}
	
\end{frame}

\begin{frame}
	\frametitle{Quelques applications : Transformation}
	\framesubtitle{Résumé automatique : Approche par graphes (Exemple : \cite{21-aries-al} - suite)}
	
	\vspace{-6pt}
	\begin{itemize}
		\item Simplification du graphe
		
		\hspace{.5cm}$noeud\_faible(v_i) = ( \sum_{(v_i, v_j) \in E} w_{ij} < \frac{1}{MImpN(v_i)} )$ 
		
		\hspace{.5cm}$arc\_faible(v_i, v_j) = ( w_{ij} < \frac{Threshold}{MImpN(v_i)})$
		
		\item Score statistique des phrases
		
		\hspace{.5cm}$ Score(s_i/ sim) = sim(s_i, C\backslash s_i) $
		$Score(s_i/ tfisf) = \sqrt{\sum\limits_{w_{ik} \in s_i} (tfisf(w_{ik}))^2}$
		
		\hspace{.5cm}$Score(s_i/ size) = \frac{1}{|s_i|}$
		$Score(s_i/ pos) = \max (\frac{1}{i}, \frac{1}{|D| - i + 1})$
		
		\hspace{.5cm}$SSF(s_i/ F) = \prod_{f_i \in F} score(s_i/f_i)$
		
		\item Score à base de graphe 
		
		\hspace{.5cm}Ex. $GC1(s_i) = SSF(s_i) + \sum\limits_{(s_i, s_j) \in E} sim(s_i, s_j) * SSF(s_j)$
		
		\item Extraction
		
		\hspace{.5cm}Ex. $ suiv_{e4}  =  \arg\min\limits_i (iord\ ssfgc(s_i) + ord\ simil(dernier{e4}, i))$ 
		
		\hspace{3cm}$ \text{ où } (dernier_{e4}, s_i) \in E $
	\end{itemize}
	
\end{frame}

\begin{frame}
	\frametitle{Quelques applications : Transformation}
	\framesubtitle{Résumé automatique : Approche linguistique}
	
	\begin{itemize}
		\item \optword{Mots de sujet} : une liste des mots pertinents au sujet, comme ``significant", ``impossible", etc.
		\begin{itemize}
			\item Ex. \cite{69-edmundson} : Bonus (mots positivement pertinents), Stigma (mots négativement pertinents)
			
			$Score_{cue}(s_i) = \sum_{w \in s_i}{cue(w)}
			\text{ où }
			cue(w) = \left\lbrace 
			\begin{array}{ll}
			b > 0 & \text{si } (w \in Bonus) \\
			\delta < 0 & \text{si } (w \in Stigma) \\
			0 & sinon 
			\end{array} 
			\right. $
		\end{itemize}
		\item \optword{Indicateurs} : Des structures qui impliquent que la phrase les contenant a une chose importante à propos du sujet
		\begin{itemize}
			\item Ex. \expword{the principal aim of this paper is to investigate ...}
		\end{itemize}
		\item \optword{Co-référence} : utilisation des anaphores ou des représentations sémantiques (Ex. Wordnet)
		\item \optword{Structure rhétorique} : utilisation de la structure rhétorique pour noter les phrases ou des syntagmes
	\end{itemize}
	
\end{frame}

\begin{frame}
	\frametitle{Quelques applications : Transformation}
	\framesubtitle{Résumé automatique : Approche linguistique (Exemple : \cite{81-paice})}
	
	\begin{itemize}
		\item Utilisation des patrons (templates) préparés manuellement
		\item\ [x] : x mots peuvent être entre ce mot et le mot précédent
		\item +y : le score est augment par y
		\item ? : ce mot est optionnel
	\end{itemize}

	\begin{figure}[!ht]
		\begin{center}
			\hgraphpage[.7\textwidth]{paice-template.pdf}
			\caption{Un exemple d'un patron simplifié \cite{81-paice}.}
			\label{fig:paice-template}
		\end{center}
	\end{figure}
	
\end{frame}

%\begin{frame}
%	\frametitle{Quelques applications : Transformation}
%	\framesubtitle{Résumé automatique : Approche linguistique (Exemple : \cite{})}
%	
%	\begin{itemize}
%		\item 
%	\end{itemize}
%	
%\end{frame}

\begin{frame}
	\frametitle{Quelques applications : Transformation}
	\framesubtitle{Résumé automatique : Approche par apprentissage automatique (AA)}
	
	\begin{itemize}
		\item \optword{Par caractéristiques}
		\begin{itemize}
			\item Réglage : régler des hyper-paramètres comme les poids des caractéristiques pour le score
			\item Classement : décider si une unité (phrase) appartient au résumé ou non
		\end{itemize}
		\item \optword{Bayesian topic models}
		\begin{itemize}
			\item Identifier les concepts principaux à partir des documents et les liens entre eux pour avoir une hiérarchie
		\end{itemize}
		\item \optword{Deep learning}
		\begin{itemize}
			\item Réglage ou Classement
			\item Génération des mots (une forme du classement)
		\end{itemize}
		\item \optword{Reinforcement learning}
		\begin{itemize}
			\item Utilisation des actions et des récompenses pour entrainer un système à générer des résumés
		\end{itemize}
	\end{itemize}
	
\end{frame}


\begin{frame}
	\frametitle{Quelques applications : Transformation}
	\framesubtitle{Résumé automatique : Approche par AA (Exemple : \cite{2020-aries})}
	
	\begin{center}
		\vgraphpage{ml2es-archi.pdf}
	\end{center}
	
\end{frame}

\begin{frame}
	\frametitle{Quelques applications : Transformation}
	\framesubtitle{Résumé automatique : Approche par AA (Exemple : \cite{06-daumeiii-marcu})}
	
	\begin{minipage}{.6\textwidth}
		\begin{itemize}
			\item $D$ : un ensemble de $K$ documents
			\item $Q$ : un ensemble de $J$ requêtes
			\item $P^G$ : un modèle de langue général de l'anglais 
			\item $P^Q$ : un modèle de langue des requêtes
			\item $P^D$ : un modèle de langue d'arrière-plan (des documents)
		\end{itemize}
	\end{minipage}
	\begin{minipage}{.38\textwidth}
		\hgraphpage{btm-daumeiii_.pdf}
	\end{minipage}
	
\end{frame}

\begin{frame}
	\frametitle{Quelques applications : Transformation}
	\framesubtitle{Résumé automatique : Approche par AA (Exemple : \cite{15-rush-al})}
	
	\begin{center}
		\hgraphpage[.35\textwidth]{2015-rush-al_.pdf}
		\hgraphpage[.35\textwidth]{2015-rush-al-exp_.pdf}
	\end{center}
	
	\begin{itemize}
		\item[(a)] Un décodeur qui cherche le prochaine mot du résumé $y_{i+1}$ sachant la phrase en entrée $x$ et des mots déjà générés pour le résumé $y_c \equiv [y_{i-c+1},\ldots, y_i]$
		\item[(b)] Un encodeur avec attention
	\end{itemize}
	
\end{frame}

\begin{frame}
	\frametitle{Quelques applications : Transformation}
	\framesubtitle{Résumé automatique : Approche par AA (Exemple : \cite{18-narayan-al})}
	
	\hgraphpage{narayan-al_.pdf}
	
\end{frame}

%===================================================================================
\section{Interaction}
%===================================================================================

\begin{frame}
	\frametitle{Quelques applications}
	\framesubtitle{Interaction}
	\begin{itemize}
		\item 
	\end{itemize}
\end{frame}

\subsection{Questions-Réponses}

\begin{frame}
	\frametitle{Quelques applications : Interaction}
	\framesubtitle{Questions-Réponses}
	\begin{itemize}
		\item 
	\end{itemize}
\end{frame}

\subsection{Systèmes de dialogue et chatbots}

\begin{frame}
	\frametitle{Quelques applications : Interaction}
	\framesubtitle{Systèmes de dialogue et chatbots}
	\begin{itemize}
		\item 
	\end{itemize}
\end{frame}


%===================================================================================
\section{Classification}
%===================================================================================

\begin{frame}
	\frametitle{Quelques applications}
	\framesubtitle{Classification}
	\begin{itemize}
		\item Filtrage de spams
		\item Identification des langues
		\item Lisibilité 
		\item Analyse des sentiments
		\item Détection de l'humeur
	\end{itemize}
\end{frame}

\subsection{Analyse des sentiments}

\begin{frame}
	\frametitle{Quelques applications : Classification}
	\framesubtitle{Analyse des sentiments}
	\begin{itemize}
		\item 
	\end{itemize}
\end{frame}

\subsection{Lisibilité}

\begin{frame}
	\frametitle{Quelques applications : Classification}
	\framesubtitle{Lisibilité}
	\begin{itemize}
		\item 
	\end{itemize}
\end{frame}


%===================================================================================
\section{Voix}
%===================================================================================

\begin{frame}
	\frametitle{Quelques applications}
	\framesubtitle{Voix}
	\begin{itemize}
		\item 
	\end{itemize}
\end{frame}

\subsection{Reconnaissance de la voix}

\begin{frame}
	\frametitle{Quelques applications : Voix}
	\framesubtitle{Reconnaissance de la voix}
	\begin{itemize}
		\item 
	\end{itemize}
\end{frame}

\subsection{Synthèse de la voix}

\begin{frame}
	\frametitle{Quelques applications : Voix}
	\framesubtitle{Synthèse de la voix}
	\begin{itemize}
		\item 
	\end{itemize}
\end{frame}


\insertbibliography{TALN10}{*}

\end{document}

